\documentclass[12pt]{article}
\usepackage[utf8]{inputenc}
\usepackage[T1]{fontenc}
\usepackage{hyperref}
\title{Shubi Language Resources for K-12 Curriculum Development}
\author{Compiled by Claude Sonnet 5.5 (claude-sonnet-5-5) for Andr\'as Kornai}
\date{}
\begin{document}
\maketitle

\section{Introduction}
Shubi (ISO 639-3 \texttt{suj}) is a Bantu language of north-western Tanzania (Kagera Region). Wikipedia gives a 1987 count of 153,000 speakers. Other names are Kishubi, Kisubi, Shuwi, Sinja, Sinya, Subi and Urushubi. Glottolog notes that a separate language, Subi (xsj), is sometimes confused with Shubi. Little material exists, and this report lists only what the searches turned up.

\section{Primary Educational Resources}
No Shubi primers, textbooks or school programmes were found.

\section{Linguistic and Reference Materials}
\begin{itemize}
\item Wikipedia, ``Shubi language''. \url{https://en.wikipedia.org/wiki/Shubi_language}
\item Glottolog entry for ISO code suj. \url{https://glottolog.org/resource/languoid/iso/suj}
\item No Shubi dictionary or grammar turned up in the searches.
\end{itemize}

\section{Religious and Community Resources}
\begin{itemize}
\item Joshua Project language page. It lists the scripture status as ``Translation Started''; the Audio Bible and Jesus Film are shown as not available. \url{https://joshuaproject.net/index.php/languages/suj}
\item Joshua Project, Shubi in Tanzania. \url{https://joshuaproject.net/people_groups/15101}
\item Wikipedia, ``Shubi people''. \url{https://en.wikipedia.org/wiki/Shubi_people}
\item African Proverbs, Sayings and Stories (Tangaza University). The search summary mentions a 2006 photocopied collection of 198 Shubi proverbs by Nkumbulwa and Tertrais, published by the Sukuma Research Committee in Mwanza. I did not confirm that this page itself holds the collection. \url{https://afriprov.tangaza.ac.ke/?p=467}
\end{itemize}

\section{Digital and Audio Resources}
No Shubi audio, video, corpus or software resources were found.

\section{Curriculum Development Context}
Shubi is not known to be taught in schools, and very little published material is available. Any native-language curriculum would need new primer, orthography and reading materials. A local proverb collection and a started Bible translation are the only leads. Kiswahili is the likely language of wider instruction in Tanzania.
\end{document}